\section{Supplementary Evaluations}
\label{app:supplemental}


\subsection{Targeted Invocation}
\label{sec:targeted-invocation-evaluation}

For a known provider, we compare NDNSF's four-message flow and targeted requests with gRPC and NSC using a 64-byte echo command and a configured 5\,ms handler wait. A call succeeds when the returned command and executor identity are validated within the configured 5\,s timeout. Five matched blocks per condition rotate arm order after setup and 10\,s warm-up; each measures 60\,s with one outstanding call. Two 1\,Gbps links have 1, 10 or 50\,ms one-way delay without loss, or 10\,ms with 1\% loss. NDNSF uses a 4096~B publication budget within a 6144~B Sync Interest Application Parameters budget; these limits are distinct from the business payload size. NDNSF uses NDNSF authorization and encryption, gRPC uses mutual TLS, and NSC validates signed packets without payload encryption; the comparison does not isolate protection overhead. Targeted calls use a prepared challenge to carry input in Request and omit ACK and Selection when eligible; the four-message flow carries input in Selection. Latencies are equal-weight averages of the five block means; first-call observations start after readiness, before warm-up, and exclude setup.

Without loss, targeted reduces successful mean latency by 38.3--48.7\% relative to the four-message flow, with all attempts completed. At 10\,ms one-way delay, their means are 80.179 and 145.649\,ms, versus 76.929\,ms for gRPC and 139.875\,ms for NSC. All 60 sampled first calls without loss succeed, but targeted has no first-call mean advantage. Four of the 20 first calls with loss time out. At 1\% loss, targeted completes 212/269 attempts (78.810\%), versus 309/355 (87.042\%) for the four-message flow, NSC's 390/440 (88.636\%) and gRPC's 3179/3179. The loss condition includes bootstrap recovery in 36 of the 212 successful targeted calls. This experiment does not attribute the loss-related failures to a specific protocol stage or establish a flight-control deadline.

\subsection{Permission Management: Methods and Detailed Results}
\label{app:measured-permission-cost}

We compare three C++ measurement configurations with identical initial permissions: eight RSA-2048 identities, each authorized for 1, 4 or 16 services. Aggregated ABE issues one DKEY containing all of an identity's permissions. Per-permission ABE uses the same NAC-ABE/OpenABE implementation but issues separate DKEYs. A batching alternative distributes AES-256 service keys in one RSA-protected bundle per identity and supports incremental grants. All three reuse signing credentials, sign their management objects, encrypt business data with AES-256 and cache unchanged material. A grant adds one permission to one identity; revocation removes that identity and rotates keys for the seven survivors. Five matched repetitions rotate their order. Functional checks verify authorized access, rejection of unauthorized or tampered data, rejection of new-generation ciphertext by a revoked identity retaining old keys, and restored access for all seven survivors.

Table~\ref{tab:permission-component-cost} reports local CPU time for parameter/key setup, material generation, signing, verification and recipient decryption. Bytes sum signed management objects and one shared public-parameter object per generation. The experimental harness packages these objects for comparison; they are not NDNSF Authority DKEY packets or measured network traffic. Business-data delivery is excluded from this management-byte column: both ABE configurations additionally publish a signed wrapped content-key object of about 819\,B for the measured 64-byte payload, whereas the RSA alternative already distributes service keys in its management bundle. Common identity bootstrap is excluded from per-configuration costs.

\begin{table}[t]
\centering
\caption{Local management costs for eight identities with 16 permissions each. CPU: mean $\pm$ 95\% confidence-interval half-width over five runs. KiB: signed management-object size, including public parameters once, excluding business data and content keys. Objects: DKEYs or RSA-protected bundles, excluding public parameters. Agg.: aggregated ABE; Sep.: separate ABE keys; Batch: batched RSA/AES.}
\label{tab:permission-component-cost}
\footnotesize
\begin{tabular}{llrrr}
\toprule
Phase & Config. & CPU (ms) & KiB & Objects \\
\midrule
Initial & Agg. & $303.10\pm2.91$ & 36.55 & 8 \\
 & Sep. & $704.58\pm16.79$ & 113.60 & 128 \\
 & Batch & $6.15\pm0.11$ & 11.27 & 8 \\
Grant & Agg. & $38.73\pm0.48$ & 4.67 & 1 \\
 & Sep. & $5.40\pm0.12$ & 0.88 & 1 \\
 & Batch & $0.71\pm0.05$ & 0.66 & 1 \\
Revoke & Agg. & $268.80\pm11.06$ & 32.09 & 7 \\
 & Sep. & $605.42\pm6.57$ & 99.51 & 112 \\
 & Batch & $5.69\pm1.34$ & 9.86 & 7 \\
\bottomrule
\end{tabular}
\end{table}

Aggregation reduces initial CPU by 57.0\% and encoded management bytes by 67.8\% relative to separate ABE DKEYs at 16 permissions, with a similar reduction during global ABE rotation. At one permission, initial CPU is 51.05 versus 51.66\,ms and both encode about 7.90\,KiB: aggregation has little room to help. Adding one permission reverses the comparison because an aggregated key replaces the complete policy, whereas the other configurations distribute a delta. Identity-wide revocation rotates the ABE generation and replaces all seven survivors' DKEYs. The RSA baseline rotates only service keys held by the withdrawn identity and updates affected survivors; in this deliberately shared access matrix all seven are affected.

The batched RSA/AES alternative has lower local management costs than aggregated ABE at every measured permission scale and phase. The component results therefore support aggregation within ABE, without establishing a general advantage over batched key distribution.

The network comparison distributes the original signed management Data from an authority-side server to eight C++ recipients, each holding only its own private keys. It uses the same permissions and retained-old-key rejection checks. Three links have 10\,ms one-way delay, 1\,Gbps capacity and no configured loss. Each recipient uses one Face, at most eight outstanding exact-name Interests, a 1\,s Interest lifetime, no retry and a 10\,s phase deadline. Fresh processes and NFD receive 10\,s routing warm-up and at least 60\,s capture. Five cryptographic repetitions cover three configurations, three scales and two certificate-cache policies: all 90 runs satisfy the prescribed access-control checks.

Management windows cover certificate and key retrieval, validation and installation; access probes run afterward. Network CPU sums recipient process time and excludes generation. Recipient processing is included in both CPU measurements, so they represent overlapping costs. IP bytes count both directions on each of the three links. NFD may cache immutable public parameters and certificates.

\begin{table}[t]
\centering
\caption{Network management costs at 16 permissions per identity with application certificate reuse. Means over five runs. Data counts retrieved management packets across recipients, including segments and public parameters; CPU sums recipient process time; IP sums traffic over three links. Material generation and access probes are excluded. Configuration labels follow Table~\ref{tab:permission-component-cost}.}
\label{tab:permission-network-cost}
\footnotesize
\begin{tabular}{llrrr}
\toprule
Phase & Config. & Data & CPU (ms) & IP (KiB) \\
\midrule
Initial & Agg. & 32 & 40.46 & 86.74 \\
 & Sep. & 136 & 125.26 & 281.72 \\
 & Batch & 8 & 36.90 & 27.25 \\
Grant & Agg. & 3 & 3.74 & 10.62 \\
 & Sep. & 1 & 4.04 & 2.16 \\
 & Batch & 1 & 2.96 & 1.71 \\
Revoke & Agg. & 29 & 12.94 & 75.80 \\
 & Sep. & 120 & 88.94 & 247.11 \\
 & Batch & 7 & 9.51 & 22.51 \\
\bottomrule
\end{tabular}
\end{table}

At sixteen permissions, eight aggregated DKEYs replace 128 separate DKEYs. Segmentation narrows the initial retrieval count to 32 versus 136 management Data packets. Aggregation also reduces initial management traffic from 281.72 to 86.74\,KiB and rotation traffic from 247.11 to 75.80\,KiB relative to separate ABE keys. Paired traffic reductions are 194.98\,KiB (95\% $t_4$ CI: 193.87--196.10) and 171.30\,KiB (170.70--171.91). Initial mean recipient management completion time falls from 328.81 to 192.85\,ms; the paired reduction is 135.96\,ms (129.17--142.75). At four permissions, initial traffic falls from 74.10 to 31.97\,KiB, but the management completion-time difference remains unresolved. At one permission, initial traffic is about 21.96\,KiB for both ABE configurations, leaving little aggregation benefit.

A grant reverses the traffic comparison: replacing the sixteen-permission policy costs an extra 8.45\,KiB (8.45--8.46), whereas separate ABE delivers one additional key. Full-policy replacement increases the cost of this single grant; cumulative costs across repeated updates remain unmeasured. Batched RSA/AES uses fewer link bytes than aggregated ABE at every measured scale and phase. This distribution experiment uses pre-generated signed material, rather than an online sequence of generation and updates. Its results do not establish recovery through the NDNSF Authority API or equivalence of every security property of the constructions.

All alternatives reuse signing identities. The certificate ablation either retains the initially validated authority certificate in the application or drops that reference between phases; NFD caching stays enabled. Reuse avoids eight certificate Interests per grant or revocation phase and saves 0.82--1.05\,ms grant recipient CPU, with all paired intervals above zero. Every grant pair has identical link traffic, and revocation traffic differences remain unresolved. Application reuse saves processing here, without demonstrated link-traffic savings.

A separate MiniNDN experiment exercises the NDNSF Controller, Provider and User update APIs with eight identities and one or sixteen initial permissions per identity. We set measurement timeouts of 10\,s for grant recovery and 20\,s for revocation recovery. Identity private keys are not shared between actors. Status polling is every 500\,ms, status validity is 24\,h and the configured grace interval is zero. All five runs at each permission scale recover protected access after the grant and restore the seven survivors after revocation and key rotation. Grant recovery ranges from 0.644 to 2.569\,s at one permission and 0.664 to 2.713\,s at sixteen; revocation recovery ranges from 2.691 to 2.819\,s and 3.348 to 7.022\,s, respectively. The Provider installs the updated revocation state for all configured services, with a maximum observed installation delay of 4.650\,s.

Across these ten runs, 102 successful calls using the pre-revocation permission state complete after the revocation update and 18 requests have no terminal result by the end of the observation window. Completion timestamps alone do not establish when an execution was authorized. We observe no execution by a revoked user starting after the updated state is installed for the exercised service; other services have no post-installation execution samples. These measurements show access recovery in this configuration, but do not establish immediate revocation or uninterrupted service. They do not compare management costs through the framework APIs across the three cryptographic alternatives.



